\section*{Chapter \ref{ch:iv1:random-walks-and-martingales} --- Random Walks and Martingales}
\markboth{\protect\omTocMark{}{Answers: Random Walks and Martingales}}{\protect\omTocMark{}{Answers: Random Walks and Martingales}}

\begin{solution}{iq:iv1:random-walks-and-martingales:1}
You end at 10 with probability $0.3$, and you take 21 steps on average.

Use optional stopping\index{optional stopping}. Your position is a martingale\index{martingale}, a process whose expected next value is its current value, so its expected value at the end equals where you started: $10p = 3$, so $p = 0.3$.

For the time, $S_t^2 - t$ is also a martingale, because one step adds 1 to the square on average. At the end, $E[S_T^2] - E[T] = 3^2$, and $E[S_T^2] = 0.3 \cdot 100 = 30$. So $E[T] = 30 - 9 = 21$, which is the general $k(n - k)$.

\iqwrong{Solving nine linear equations, which is correct but slow. Using the two martingales saves that work, and they extend to harder questions.}
\end{solution}

\begin{solution}{iq:iv1:random-walks-and-martingales:2}
You reach \$10 with probability about $0.884$, and it takes about $19.2$ bets on average.

With an uneven coin, $r^{S_t}$ is a martingale\index{exponential martingale}, where $r = q/p = 2/3$: over one step, $p r + q/r = q + p = 1$. Starting at 5, and with $P$ the chance of reaching 10, optional stopping gives
\[
r^5 = P r^{10} + (1 - P) \quad\text{so}\quad P = \frac{1 - r^5}{1 - r^{10}} = \frac{1}{1 + r^5} \approx 0.8836.
\]
For the time, use Wald's identity\index{Wald's identity}. Your average gain per bet is $0.2$, and your average final gain is $10P - 5 \approx 3.84$, so the number of bets is $3.84/0.2 \approx 19.2$.

\iqwrong{Assuming that the 60\% win rate means a 60\% chance of doubling. The edge compounds over many bets, so the chance is much higher.}
\iqwhy{A small edge gets close to certainty once the target is many bets away. This is what makes position sizing, and not forecasting alone, decide risk.}
\end{solution}

\enlargethispage{-3\baselineskip}%
\begin{solution}{iq:iv1:random-walks-and-martingales:3}
You win the \$1 in 90\% of sessions. The average profit is 0, and a session takes 9 bets on average.

A fair game stays fair when stopped\index{optional stopping}, as long as the stopping is bounded like this. So $0 = P \cdot 1 + (1 - P)(-9)$, which gives $P = 0.9$. With $a = 1$ and $b = 9$ the general result for the time is $ab = 9$.

\iqwrong{Thinking that a 90\% win rate shows a good strategy. The rare loss is exactly large enough to cancel all the wins. A high hit rate says nothing about whether a game is fair, or about what it is worth.}
\iqwhy{Selling insurance looks like this: many small wins, one large loss. A track record with a 90\% hit rate may be nothing but the walk between two stops.}
\end{solution}

\begin{solution}{iq:iv1:random-walks-and-martingales:4}
Back at 0 after 10 steps has probability $252/1024 \approx 0.246$. After 100 steps, the average distance is about $7.96$, and $\sqrt{2N/\pi}$, with $N = 100$ steps, gives $7.98$.

To be at 0 after $2n$ steps, exactly $n$ of them must go up, so $u_{2n} = \binom{2n}{n}/4^n$.

For the distance, there is an exact formula\index{mean absolute displacement}: $E|S_{2n}| = 2n\,u_{2n}$, here $100 \cdot 0.0796$. The approximation comes from a normal with variance $N$, whose mean absolute value is $\sigma\sqrt{2/\pi}$.

\iqwrong{Answering 10, the standard deviation. The average distance is smaller, by a factor of $\sqrt{2/\pi} \approx 0.8$.}
\end{solution}

\enlargethispage{-3\baselineskip}%
\begin{solution}{iq:iv1:random-walks-and-martingales:5}
It takes 50 steps on average.

Use Wald's identity\index{Wald's identity}: if a walk with independent steps of mean $\mu$ stops at a stopping time $T$ with finite mean, then $E[S_T] = \mu E[T]$. Here each step has mean $0.2$, and the walk stops exactly at $+10$, because it moves one place at a time. So $E[T] = 10/0.2 = 50$.

A second route checks the identity. Let $x$ be the mean time to go up one level. The first step goes up with $0.6$. Otherwise you must climb two levels: $x = 1 + 0.4 \cdot 2x$, so $x = 5$, and ten levels take 50.

\iqwrong{Using $10/0.6$, by counting only the up steps. The down steps have to be recovered, and it is the net drift that sets the rate.}
\iqfollow{And as the drift goes to zero?}{The time grows as $10/\mu$ without limit: at zero drift it is infinite, even though the walk still reaches $+10$ with certainty.}
\end{solution}

\begin{solution}{iq:iv1:random-walks-and-martingales:6}
It reaches $+10$ with probability $(2/3)^{10} \approx 0.0173$, and its highest point averages 2.

Let $x$ be the chance of ever going up one level. Using first-step analysis\index{first-step analysis}, you go up at once, or go down and must rise two levels: $x = 0.4 + 0.6 x^2$. The roots are $2/3$ and 1. The walk drifts downward, so it is not certain to rise, and the small root is right. To reach $+10$, you have to rise one level ten times in a row, and each level is a new, independent problem, so the chance is $x^{10}$.

So the highest point $M$ has $P(M \ge k) = (2/3)^k$, a geometric distribution\index{maximum of a negative-drift walk}, and
\[
E[M] = \sum_{k \ge 1} (2/3)^k = \frac{2/3}{1 - 2/3} = 2.
\]
\iqwrong{Taking the root $x = 1$ because the equation allows it. With a negative drift, the walk heads off to minus infinity, so a limit on the maximum must exist.}
\iqwhy{This gives the distribution of the deepest loss below your starting capital for a positive-edge strategy on a fixed bet: look at it from the other side, and it is the highest point of a negative walk. It is not the maximum drawdown, which measures from each new peak and keeps growing the longer you trade.}
\end{solution}

\begin{solution}{iq:iv1:random-walks-and-martingales:7}
You end \$1 ahead with probability $1 - 2^{-10} \approx 0.999$. The expected profit is 0, and a session has about $1.998$ bets on average.

The bets $1, 2, 4, \dots, 512$ add up to \$1\,023, so you can make 10 bets. You lose them all with probability $2^{-10}$, and then you lose all \$1\,023. Otherwise you win \$1, because any win covers all the earlier losses plus one. So $E[\text{profit}] = (1 - 2^{-10}) - 2^{-10} \cdot 1023 = 0$.

The number of bets is capped at 10, and adding up $P(\text{bet } k \text{ happens}) = 2^{-(k - 1)}$ gives $2 - 2^{-9}$.

\iqwrong{``Nearly always a profit, so this must be good.'' By optional stopping\index{optional stopping}, no stopping rule makes a fair game profitable when capital is finite. The tiny loss chance carries a loss of a thousand times the win.}
\iqfollow{And with unlimited capital?}{You always win \$1, in about 2 bets, but the capital you may need has an infinite mean. The system needs what no one has.}
\end{solution}

\begin{solution}{iq:iv1:random-walks-and-martingales:8}
It does not contradict the claim, because the theorem needs conditions that fail here: the stopping time has infinite mean, and the walk is unbounded below. The median stopping time is 1, while the mean is infinite.

Optional stopping\index{optional stopping} holds for a bounded time, or for a bounded martingale. Here neither holds. At any fixed horizon the game is fair: in the paths that have not stopped, the average position is well below 0, and this balances the paths that stopped at $+1$.

The time $T$ is 1 with probability $1/2$, which fixes the median. By reflection\index{reflection principle}, $P(T > 2n) = u_{2n} = \binom{2n}{n}/4^n \approx 1/\sqrt{\pi n}$. For example, $P(T \le 100) \approx 0.92$. The tail falls too slowly for the sum $\sum_n P(T > n)$ to converge, so $E[T] = \infty$. You would have to survive arbitrarily large losses in the meantime.

\iqwrong{Concluding that the walk cannot be certain to reach $+1$, because the game is fair. It is certain to reach it, and fairness at every fixed time does not prevent that.}
\iqwhy{A strategy that holds losers until they recover is this walk. It shows a profit every time it stops, and the cost is hidden in the unbounded time and the unbounded loss you have to carry.}
\end{solution}

\begin{solution}{iq:iv1:random-walks-and-martingales:9}
From 0 to 10 takes 100 steps on average, and from 4 to 10 takes 84.

Use $S_t^2 - t$\index{quadratic martingale}. On the interior, one step adds 1 to $S^2$ on average, so this stays fair. At 0, the forced step takes $S^2$ from 0 to 1, which also adds exactly 1, so the same process stays fair at the reflecting point. Stopping at 10 gives $E[S_T^2] - E[T] = k^2$, and so $E[T] = 100 - k^2$.

For $k = 0$ that is 100, and for $k = 4$ it is 84.

\iqwrong{Using 24 steps from 4, the $k(n - k)$ of a walk with an exit at 0. That walk would stop at 0, while this one bounces back and carries on.}
\iqfollow{In general?}{Reaching a level $n$ from a reflecting wall takes $n^2$ steps. Time grows as the square of the distance, the signature of a diffusion.}
\end{solution}

\enlargethispage{-4\baselineskip}%
\begin{solution}{iq:iv1:random-walks-and-martingales:10}
The probability is about $0.928$, and it takes about $21.4$ steps on average.

The steps at which the walk stays put change nothing about where it goes next, so ignore them. The moving steps go up with probability $0.5/0.8 = 5/8$. So $r = 3/5$, and the exit formula\index{exponential martingale} gives
\[
P = \frac{1}{1 + r^5} = \frac{1}{1 + 0.6^5} \approx 0.9279.
\]
The time does change. By Wald's identity\index{Wald's identity}, the final position averages $5P - 5(1 - P) \approx 4.279$. The mean step, including the zero, is $0.5 - 0.3 = 0.2$. So the time is $4.279/0.2 \approx 21.4$.

\iqwrong{Leaving the stay-put steps out of the duration, or letting them change the exit probability. They are irrelevant for where you exit, but not for how long you take.}
\iqwhy{Quote-update data often has many unchanged steps. Filtering those out changes the clock but not the direction, and it is cleaner to separate these two effects.}
\end{solution}

\begin{solution}{iq:iv1:random-walks-and-martingales:11}
The target is hit first with probability about $0.868$, the trade lasts about 160 ticks, and it earns about $16.03$ ticks on average. A fair walk would give a chance of $1/3$ and an expected profit of 0.

With $r = 0.45/0.55 = 9/11$, the exit formula\index{exponential martingale} from 10 ticks above the stop and with 30 between stop and target is
\[
P = \frac{1 - r^{10}}{1 - r^{30}} \approx 0.8677.
\]
The expected profit is $20P - 10(1 - P) = 30P - 10 \approx 16.03$ ticks. Each tick earns $0.1$ on average, so Wald's identity\index{Wald's identity} gives $16.03/0.1 \approx 160$ ticks.

With a fair walk, $P = 10/30 = 1/3$, and the profit is zero whichever stop and target you choose. Stops and targets change the shape of the outcome, not its mean. A trader's edge comes from the tick probabilities, and these exits only harvest it.

\iqwrong{``A target twice the stop, so I need only a third of the trades to win.'' That is a statement about break-even, not an edge. It can be true of a fair walk, and it earns nothing.}
\iqfollow{How long does a fair walk take?}{From the martingale $S_t^2 - t$, the fair-walk time is the product of the distances to the two exits: $10 \cdot 20 = 200$ ticks on average. Longer than with the drift, because you get no help moving towards the target.}
\end{solution}

\begin{solution}{iq:iv1:random-walks-and-martingales:12}
About $1.9 \times 10^{-9}$ for \$1 bets, $0.134$ for \$10, and $0.669$ for \$50.

Scale the walk\index{risk of ruin}. With stakes of $s$, your capital is $100/s$ bets, and the chance of ever losing them is $r^{100/s}$, with $r = q/p = 0.45/0.55 = 9/11$: it is the chance of a walk with a positive drift ever falling that many levels. So you need $r^{100}$, $r^{10}$ and $r^2$.

Note what does not change: each bet has the same edge in percent. Expected profit per dollar bet is the same, but ruin rises from impossible to likely as the stake grows, because your capital measured in bets is what protects you.

\iqwrong{Thinking of ruin as depending on the edge alone. With an edge you still go broke if your stake is too large, and a fair coin goes broke for certain at any stake.}
\iqwhy{Risk of ruin decays exponentially in capital measured in bets. Hence stake size matters more than almost anything else for survival: doubling the bet replaces the chance of ruin by its square root, which is far larger.}
\end{solution}

\begin{solution}{iq:iv1:random-walks-and-martingales:13}
It stays strictly above 0 for 20 steps with probability $\tfrac12 u_{20} \approx 0.0881$. It first returns at step 2, 4 and 6 with probability $1/2$, $1/8$ and $1/16$.

Here $u_{2n} = \binom{2n}{n}/4^n$ is the chance of being at 0 at time $2n$. By the ballot theorem\index{ballot theorem}\index{reflection principle}, the paths that stay positive for $2n$ steps number exactly half of those ending at 0, and that gives $\tfrac12 u_{2n}$. So $u_{20} \approx 0.176$ is the chance of no return at all in 20 steps, in either direction.

For the first return, the same ideas give $f_{2n} = u_{2n}/(2n - 1)$\index{first return time}. So $f_2 = 1/2$, $f_4 = (3/8)/3 = 1/8$ and $f_6 = (5/16)/5 = 1/16$. Listing all the paths confirms this. The walk returns with certainty, but the mean of the time is infinite.

\iqwrong{Using $2^{-20}$, which treats ``stays above 0'' as needing every step to go up. Most paths that stay above 0 go up and down a lot.}
\iqwhy{A fair process can stay on one side for a long time with quite a high probability. A run of good P\&L that lasts 20 periods is only about 1 in 11.}
\end{solution}

\begin{solution}{iq:iv1:random-walks-and-martingales:14}
The probability is $(1/3)/(8/9) = 3/8$. The answer relies only on the fact that a correct belief is a martingale\index{belief as a martingale}.

If your belief is the correct conditional probability, then its expected next value is its current value: otherwise you would already believe something else. So stop the first time it reaches $8/9$. The odds move in factors of 2, so it lands exactly on $8/9$. If it never gets there, it heads to 0 in the end, because it converges to the truth. Optional stopping\index{optional stopping} gives
\[
\tfrac13 = P \cdot \tfrac89 + (1 - P) \cdot 0, \quad\text{so}\quad P = \tfrac38.
\]
Neither the news's reliability nor the number of days appears, because a correct belief moves so that this holds in any case. Odds running from $1 : 2$ up to $8 : 1$ make it a clean walk, and a simulation with a true state drawn from the prior gives the same result.

\iqwrong{Giving $8/9$, or thinking that because the event is unlikely, the belief cannot rise. A belief at $1/3$ will often rise and then fall back, because an event at $1/3$ fails two times in three.}
\iqwhy{A price of a binary contract is such a belief. At 30 cents, it can touch 90 at most a third of the time, whatever happens.}
\end{solution}

\begin{solution}{iq:iv1:random-walks-and-martingales:15}
On average they take 16 steps.

Track only their separation\index{state lumping}. It starts at 4 on a ring of 12. At each step, each token moves $\pm 1$, so the separation changes by $-2$ with probability $1/4$, $+2$ with probability $1/4$, and 0 with probability $1/2$. Since separations stay even, measure in units of 2: the separation lives on a ring of 6 positions, starts 2 from the meeting point, and is a walk that stays put half of the time.

A walk on a ring from distance 2 of 6 is the same as a gambler's ruin\index{gambler's ruin} on a line from 2 between 0 and 6, so it takes $2 \cdot 4 = 8$ moves. Half of the time nothing moves, so that doubles to 16 steps.

\iqwrong{Using $4 \cdot 8 = 32$, by treating the separation as one simple walk on 12 positions. The two tokens move together, which makes the separation move two places at a time, and half the time not at all.}
\iqfollow{What if they begin 3 apart?}{They never meet, because the parity of the separation never changes when both move.}
\end{solution}

\begin{solution}{iq:iv1:random-walks-and-martingales:16}
The last visit to 0 is most likely at the very start or the very end: at time 0 and at time 20, each with probability about $0.176$, and time 10 is the least likely, at about $0.061$. In the limit, the last visit falls in the first quarter with probability $1/3$.

The chance that the last visit is at time $2k$ is the chance of being at 0 then, times the chance of no return during the remaining time, and the second is also a $u$ value: $P = u_{2k}\, u_{2n - 2k}$. So the distribution is symmetric, and it has large values at both ends, because $u$ falls with time. For $n = 10$: $u_{20} \approx 0.176$, at both ends, and $u_{10}^2 \approx 0.061$ in the middle.

In the limit this becomes the arcsine law\index{arcsine law}: the fraction $x$ of the run at which the last visit occurs has
\[
P(\text{last visit before } x) = \frac{2}{\pi} \arcsin\sqrt{x},
\]
which gives $1/3$ at $x = 1/4$. So with chance $2/3$ the last quarter or the first quarter contains it.

\iqwrong{Thinking of the last visit as coming near the middle of the run, or that a fair walk should spend about half the time on each side. In truth, a fair walk typically stays on one side for most of its life.}
\iqwhy{A flat strategy with a long winning stretch or a long losing stretch is exactly what chance produces. The arcsine law warns you not to infer regime changes from long runs of one sign in cumulative P\&L.}
\end{solution}

\begin{solution}{iq:iv1:random-walks-and-martingales:17}
The first player runs out after 300 rounds on average, and one player holds all 60 after 1\,100 rounds on average. The general formulas are $3abc/(a + b + c)$ and $ab + bc + ca$.

Find quantities that go down steadily\index{martingale from a product}. When player $x$ gives a coin to player $y$, with $z$ the third player, the product $abc$ changes by $z(x - y - 1)$. Over the six ordered pairs the $x - y$ parts cancel, and the average change is $-(a + b + c)/3$, a constant, because the total stays fixed. So $abc + t(a + b + c)/3$ is a martingale. It ends at the moment someone has no coins, when the product is 0, so
\[
abc = E[T] \cdot \frac{a + b + c}{3}, \quad\text{so}\quad E[T] = \frac{3abc}{a + b + c} = \frac{3 \cdot 6000}{60} = 300.
\]
For the second part, use $ab + bc + ca$. Any move changes it by $x - y - 1$, and the average is $-1$, for three players and for two. So it falls by one per round on average, and it reaches 0 when one player holds everything. So $E[T] = 200 + 600 + 300 = 1100$.

Solving the full linear system confirms both formulas for small starts, such as 3 and 11 from $(1, 2, 3)$.

\iqwrong{Trying to write down the recursion for the state. There are many states, and the point of the question is that a good quantity gives the answer in one line.}
\iqfollow{And with two players, $a$ and $b$?}{The product $ab$ falls by one per round, which recovers the $ab$ of the gambler's ruin.}
\end{solution}

\begin{solution}{iq:iv1:random-walks-and-martingales:18}
The upper bound is $2/3$, the lower bound from the long run is $1/4$, and a direct calculation gives about $0.44$.

Martingale: if red is a fraction $f = r/(r + b)$, the next fraction has expected value
\[
f \cdot \frac{r + 1}{r + b + 1} + (1 - f) \cdot \frac{r}{r + b + 1} = \frac{r^2 + r + br}{(r + b)(r + b + 1)} = f.
\]
Doob's maximal inequality\index{Doob's maximal inequality} for a non-negative martingale gives $P(\sup \ge c) \le f_0/c$, here $(1/2)/(3/4) = 2/3$. This is just optional stopping, with the fraction at least $c$ on that event.

Lower bound: the fraction converges\index{P\'olya urn}, and its limit is uniform on $[0, 1]$. So it ends above $3/4$ with chance $1/4$, and that counts among the times it reached it. A simple sharper floor comes from two red draws in a row, which already give $3/4$: $\tfrac12 \cdot \tfrac23 = 1/3$.

A dynamic programme over the urn's counts up to 4\,000 balls gives about $0.44$, which a simulation confirms. The fraction often passes $3/4$ early, when few balls are in the urn, and later falls back. That is why the limit's chance is well below the chance of ever touching.

\iqwrong{Using the bound as the answer. Doob's bound is attained only when the path hits the level exactly and otherwise ends at 0, and here the fraction cannot fall to 0.}
\end{solution}

\begin{solution}{iq:iv1:random-walks-and-martingales:19}
The average is about $7.50$. The continuous value $\sqrt{2n/\pi}$ is $7.98$, and the correction of $\tfrac12$ gives $7.48$.

The reflection principle\index{reflection principle} gives, for $a \ge 1$,
\[
P(M \ge a) = P(S_n \ge a) + P(S_n > a),
\]
because each path that reaches $a$ and ends below it can be reflected after its first touch of $a$ to one ending above it. Then $E[M] = \sum_{a \ge 1} P(M \ge a)$, and adding the binomial tails gives $7.50$, which a full enumeration matches on shorter walks. Paths at $a$ itself are counted once, not twice, and this creates the difference from the Brownian formula.

In the continuous limit\index{maximum of Brownian motion}, $M$ has the same distribution as $|W_n|$, whose mean is $\sqrt{2n/\pi}$. The discrete walk takes steps of a whole unit, so it falls about half a step short of the limit, and the correction of $\tfrac12$ comes from that.

\iqwrong{Giving 10, or doubling $E[S_n^+]$ with no correction. The walk's standard deviation is not the mean of its maximum. In short walks and for barriers, the half-step correction is material.}
\iqwhy{A discretely monitored barrier sits effectively further away than the same continuous barrier. The usual correction shifts the barrier by a fraction of one step's standard deviation, which reflects this same fact.}
\end{solution}

\begin{solution}{iq:iv1:random-walks-and-martingales:20}
The variance is 1\,200, so the standard deviation is about $34.6$ steps.

Each step has mean $\mu = 0.2$ and variance $\sigma^2 = 1 - \mu^2 = 0.96$. The process $(S_t - \mu t)^2 - \sigma^2 t$ is a martingale\index{Wald's second identity}. Stop it at $T$, when $S_T = 10$ exactly:
\[
E[(10 - \mu T)^2] = \sigma^2 E[T].
\]
Since $E[T] = 10/\mu$, the left side is $\mu^2\,\mathrm{Var}(T)$. So
\[
\mathrm{Var}(T) = \frac{\sigma^2 E[T]}{\mu^2} = \frac{10\,\sigma^2}{\mu^3} = \frac{10 \cdot 0.96}{0.008} = 1200.
\]
A simulation of 20\,000 walks gives the same result. As you would expect, the variance grows linearly with the distance, because each level is a separate independent wait.

\iqwrong{Using a variance of $0.96 \cdot 50 = 48$, the variance of the position at a fixed time. The wait is much more variable, since the variance is divided by $\mu^3$, and a small drift gives a huge spread.}
\iqwhy{For a strategy with a small edge, the time to reach a target has a very wide distribution: here a standard deviation of 70\% of the mean. You should expect long periods with no progress, even when the edge is real.}
\end{solution}

\begin{solution}{iq:iv1:random-walks-and-martingales:21}
It returns with probability one in one and two dimensions, and with probability about $0.34$ in three dimensions. This is P\'olya's theorem\index{P\'olya's theorem}.

The number of visits to the start is a geometric count: with return chance $R$, its mean is $U = 1/(1 - R)$. So the walk returns with certainty exactly when $U = \sum_n u_n$ diverges.

In one dimension $u_{2n} = \binom{2n}{n}/4^n \sim 1/\sqrt{\pi n}$, and the sum diverges. In two dimensions, rotating the lattice by 45 degrees turns the walk into two independent one-dimensional walks, so $u_{2n}$ is the square of the 1-d value, about $1/(\pi n)$, and the sum still diverges, only just. In three dimensions, $u_{2n}$ falls as about $n^{-3/2}$, and the sum converges.

The value uses the Fourier form of $U$, the average of $1/(1 - \tfrac13(\cos a + \cos b + \cos c))$ over the cube of angles, which comes to about $1.516$. So $R = 1 - 1/1.516 \approx 0.34$. A direct simulation gives the same result.

\iqwrong{Thinking that a walk must come back because it has infinite time. In a high enough dimension, there are too many places to go.}
\iqfollow{What are the odds in four dimensions?}{Lower still, about $0.19$ on the same reasoning. The return probability falls steadily with dimension.}
\end{solution}

\begin{solution}{iq:iv1:random-walks-and-martingales:22}
You stop at the top with probability about $0.326$, in about $25.9$ steps. One-third is close, but the upper exit can overshoot to 11.

The walk has mean step 0, and it moves down one place at a time, so it lands exactly on $-5$. Going up, though, it lands on 10 or 11. Let $a$ and $b$ be the chances of stopping at 10 and at 11. Optional stopping\index{optional stopping} says $10a + 11b - 5(1 - a - b) = 0$, which gives $15a + 16b = 5$. This has two unknowns, so you need one more fact, and a linear system on the positions supplies it, with $a \approx 0.218$ and $b \approx 0.108$. Their total is $0.326$.

For the time, each step has variance 2 (it is $\tfrac13 \cdot 4 + \tfrac23 \cdot 1$), so $S_t^2 - 2t$ is a martingale. Then $E[T] = \tfrac12 E[S_T^2] = \tfrac12(100a + 121b + 25(1 - a - b)) \approx 25.9$.

If the walk landed exactly on 10, as a walk with unit steps does, you would get $5/15 = 1/3$; with this walk's step variance of 2, the time would be $5 \cdot 10/2 = 25$. (Unit steps, with variance 1, take 50.) Overshooting means the top is on average a bit further than 10. So you exit there a little less often.

\iqwrong{Ignoring overshoot because the answer looks close. A one-third answer is off by only about 2\%, but it is wrong, and with larger jumps it can be wrong by much more. Wald's identity needs the position at the stop, and this is not the barrier when jumps can pass it.}
\iqwhy{A stop order filled through a gap is an overshoot. Your actual stop is not the barrier, and a jump in the price is what breaks the textbook calculation.}
\end{solution}

\enlargethispage{-1\baselineskip}%
\begin{solution}{iq:iv1:random-walks-and-martingales:23}
The variance is $\tfrac23 n^2(n^2 - 1)$, which is 6\,600 for $n = 10$, so the standard deviation is about 81 against a mean of 100.

Use a fourth-power martingale\index{fourth-moment martingale}. Over one step, $E[(S + 1)^4 + (S - 1)^4]/2 = S^4 + 6S^2 + 1$. Correcting for this, you can check that
\[
M_t = S_t^4 - 6tS_t^2 + 3t^2 + 2t
\]
stays fair over one step. Stop it at $T$, with $|S_T| = n$:
\[
0 = n^4 - 6n^2 E[T] + 3E[T^2] + 2E[T].
\]
Since $E[T] = n^2$, this gives $3E[T^2] = 5n^4 - 2n^2$. So $E[T^2] = \tfrac13(5n^4 - 2n^2)$, and subtracting $n^4$,
\[
\mathrm{Var}(T) = \frac{2n^4 - 2n^2}{3} = \frac23 n^2(n^2 - 1).
\]
Check: at $n = 1$, you exit after one step for certain, and the formula gives 0. Second-moment first-step equations give the same value for $n$ up to 10, and a simulation matches.

\iqwrong{Assuming that the relative spread shrinks as $n$ grows. The ratio of standard deviation to mean tends to $\sqrt{2/3}$, so this stays a very variable wait.}
\iqfollow{Why is the ratio fixed?}{By scaling. A walk at scale $n$ behaves like Brownian motion on $[-1, 1]$, with time rescaled by $n^2$, so all of the moments of $T/n^2$ have a limit.}
\end{solution}

\begin{solution}{iq:iv1:random-walks-and-martingales:24}
It visits $+k$ exactly once on average, for every $k$, including 100.

For the proof, think about how the walk reaches $k$. Step first to $+1$, with chance $1/2$. Then reach $k$ before 0, by gambler's ruin\index{gambler's ruin}, with chance $1/k$. Now count the visits to $k$ before 0. From $k$ a step goes one way or the other. Going up, it always returns. Going down, it reaches 0 before returning with chance $1/k$. So each visit has an escape chance of $\tfrac12 \cdot \tfrac1k$, and the number of visits is geometric, with mean $2k$. Overall:
\[
\tfrac12 \cdot \tfrac1k \cdot 2k = 1.
\]
This is what is expected from the stationary measure\index{stationary measure}. The counting measure is stationary for the walk, and between returns to any state the expected number of visits to another state is the ratio of stationary masses, here $1$. On a finite ring this holds exactly, for all $k$, as solving the linear system confirms.

It is not absurd, since the walk reaches $+100$ only about once in 200 excursions. But when it does, it visits about 200 times. Rare excursions make up for it, and because every level contributes $1$, the total time is infinite.

\iqwrong{Saying that the far level must be visited much less often, because most returns occur quickly. Most excursions never get there, but the few that do stay around for a long time.}
\iqwhy{Rare and long excursions dominate means. A mean over many short excursions misses them in a small sample, and it may even be undefined. That is a warning about what a small backtest tells you about a heavy-tailed process.}
\end{solution}

\begin{solution}{iq:iv1:random-walks-and-martingales:25}
At \$1 a bet your chance is about $6.2 \times 10^{-14}$, and betting boldly gives $0.16$, which is the best possible. At a win rate of $0.6$ the small bets are best, with a chance near 1. At $0.5$ every plan gives $1/4$.

Betting \$1 is a gambler's ruin\index{gambler's ruin} with $r = q/p = 3/2$, starting at 25 with a target of 100:
\[
P = \frac{1.5^{25} - 1}{1.5^{100} - 1} \approx 6.2 \times 10^{-14}.
\]
Bold play\index{bold play} stakes all \$25, then \$50: two wins in a row, which has chance $0.4^2 = 0.16$. Value iteration over all whole-dollar stakes gives $0.16$ as the best value too.

Why? Each bet at $p < 1/2$ costs you in expectation, in proportion to the amount staked. Small bets force you to stake a large total before you can reach the target, and each dollar you stake loses money on average. Bold play puts up the smallest total stake that could reach the target.

At $p = 0.6$ the logic reverses: each bet earns, so stake a lot in total. The chance with small bets is $(1 - r^{25})/(1 - r^{100})$ with $r = 2/3$, which is about $0.99996$, compared with a bold $0.36$. At $p = 0.5$, your wealth is a martingale, so optional stopping says the chance is $25/100$ under any plan.

\iqwrong{Saying that small bets reduce risk, so they must be safer here. Small bets reduce the variance, and when you need to reach a target against a negative edge, variance is what you are buying.}
\iqwhy{This applies to a trader who is behind on a target with no real edge: the only way to reach it is by raising size, and this is exactly why risk limits exist. Conversely, a positive edge is destroyed by betting as if it were a lottery.}
\end{solution}
